\documentclass[UTF8,a4paper,12pt]{ctexart}

\usepackage{geometry}
\usepackage{setspace}
\usepackage{enumitem}
\usepackage{hyperref}
\usepackage{fancyhdr}

\geometry{left=2.5cm,right=2.5cm,top=2.6cm,bottom=2.6cm}
\setstretch{1.18}
\setlist[itemize]{leftmargin=2em}
\pagestyle{fancy}
\fancyhf{}
\fancyhead[C]{104754252295 王曼 软件学院 软件工程专业 研究生一年级}
\fancyfoot[C]{\thepage}
\setlength{\headheight}{15pt}

\title{我的 Vibe Research 实践：案例分析与经验教训}
\author{姓名：\underline{\hspace{4cm}} \quad 学号：\underline{\hspace{4cm}}}
\date{}

\begin{document}

\maketitle

\section{引言：从“让模型涨点”到“理解错误从哪里来”}

本学期我在机器学习课程之外，围绕开放词汇三维实例分割做了一段比较完整的 Vibe Research 实践。我的任务大致是：在已有 Mask3D 三维分割结果的基础上，尝试利用 YOLO-World 和 SAM 这类二维视觉模型，补全 Mask3D 漏掉的物体或不完整的物体区域。简单说，就是希望把二维开放词汇识别能力引入三维场景，让系统能发现更多 ScanNet200 场景中的物体。

刚开始时，我对 Vibe Research 的理解比较朴素：让 AI 帮我想实验方向、写代码、跑结果，然后看 AP 有没有提升。如果 AP 提高，就说明方向对；如果 AP 降低，就换一个想法继续试。这个过程确实很快，AI 能帮我生成大量脚本，也能快速总结实验现象。但是随着实验越来越多，我逐渐意识到，单纯依赖“感觉”和最终 AP 是非常危险的。AP 是最后的指标，但它不能告诉我错误到底来自二维检测、SAM 掩码、三维投影、Mask3D 原始预测，还是候选融合策略。

因此，我这次实践中最大的转变不是某一个具体模型带来了提升，而是研究方式发生了变化：从“不断试新模块追 AP”，转向“建立中间诊断，理解每一个候选为什么应该被接受或拒绝”。下面我结合几个亲身经历过的案例，说明这次 Vibe Research 中哪些尝试有价值，哪些做法后来证明是不对的，以及我从中得到的经验教训。

\section{案例一：盲目追 AP 是最早的错误}

最早阶段，我尝试过很多直觉上合理的方向。例如，用 Alpha-CLIP 对候选重新打分；尝试更换或补充开放词汇检测器；调整 SAM 多掩码选择；加入视角质量门控；改变候选来源的融合比例；尝试自适应内部种子；以及对三维点做简单连通分量清理。这些想法单独看都不荒谬，甚至很多都符合常识：二维模型更强，三维结果似乎就应该更好；掩码更干净，候选就应该更可靠；视角质量更高，投影误差就应该更小。

我印象比较深的是，最开始我非常容易被“这个想法听起来合理”所吸引。例如，Alpha-CLIP 重打分听起来可以让候选更符合文本语义；SAM 多掩码选择听起来可以解决单个掩码不准的问题；视角质量门控听起来可以过滤掉模糊或遮挡严重的视角；自适应内部种子听起来可以减少背景点污染。每一次提出新想法时，我都很容易觉得“这应该会有提升”。AI 也会很快帮我把这些想法转成脚本和参数组合，于是实验迭代速度变得很快。

但是实际结果并不稳定。有些实验在少量场景上看起来不错，扩到更多场景后就失效；有些策略提升了某些类别，却伤害了另一些类别；还有些修改让候选数量增加，但最终 AP 并没有改善。例如，一个更宽松的二维候选策略可能召回更多漏检物体，但同时也会把墙面、地面、柜体边缘等背景区域带进来；一个更严格的过滤策略可能让候选看起来更干净，但又会丢掉小物体或遮挡物体。单看最终 AP，我只能知道“这次分数高了或低了”，却不知道具体原因。

这个阶段还有一个很典型的问题：我经常把“候选数量变多”误以为“召回能力变强”。后来我才意识到，三维实例分割里的候选不是越多越好。一个错误候选如果和已有 Mask3D 预测重叠，可能会产生重复实例；如果跨越两个物体，可能会拉低 IoU；如果只是背景碎片，可能会增加假阳性。也就是说，候选数量增加只是表面现象，关键在于候选是否真的对应一个物体，以及它是否补充了已有三维预测没有覆盖的部分。

这一阶段最明显的问题是，我把最终 AP 当成了唯一判断标准，却没有足够的中间证据解释为什么 AP 变化。当 AP 下降时，我只能猜测是不是二维框不准、SAM 掩码不好、投影深度不一致、候选融合太激进，或者 Mask3D 本身已经覆盖了该物体。当 AP 上升时，我也不知道提升来自哪类物体、哪种候选、哪条规则。这样做实验虽然快，但每次实验留下的知识很少。

后来我反思，这种做法其实很像“用最终分数撞运气”。在三维开放词汇分割任务中，一个候选最终是否有用，受到许多环节影响：二维检测框是否正确，SAM 掩码是否覆盖了真实物体，回投影深度是否一致，多视角是否看到同一个物体，已有 Mask3D 是否已经解释了这个区域，候选是否跨越了物体边界。这些问题如果只用一个最终 AP 来判断，就很难定位错误。

这一阶段给我的第一个教训是：Vibe Research 不能只靠快速试错。AI 可以帮助我快速产生想法和代码，但如果没有诊断指标，快速试错很容易变成快速迷路。研究中真正重要的不是“又试了一个新方法”，而是每次试验后能不能更清楚地知道错误发生在哪里。后来我开始要求每个实验都尽量留下可解释的中间结果，例如候选数量、候选来源、被过滤原因、与已有 Mask3D 的重叠、点级连通分量、冲突点比例等。只有这样，实验失败才不是单纯浪费，而能变成下一轮规则修改的依据。

\section{案例二：点级证据图有帮助，但不能解决边界问题}

在意识到直接融合候选不可靠之后，我开始构建点级证据图。大致思路是：先用 YOLO-World 和 SAM 生成二维观测，再把二维掩码回投到三维点云中。不同视角看到的同一类物体，可能对应同一个三维实例；如果两个观测在三维空间中互相支持，就把它们连接起来；如果它们深度冲突或类别冲突，就标记为冲突关系。这样可以把“二维观测是否属于同一个三维物体”这个问题显式建模。

这个方向比一开始盲目融合要好很多。因为我不再直接把所有候选塞进最终预测，而是先导出证据图候选，观察支持边、弱边、冲突边、候选数量、未分配观测等中间结果。通过这些诊断，我能看到某些候选为什么被合并，某些观测为什么被拒绝，也能发现一些规则错误。例如，如果两个二维掩码来自不同视角，并且回投影后共享较多三维点，那么它们可能在支持同一个物体；如果两个掩码在同一帧中深度层次明显不同，那么它们更可能是遮挡关系或不同物体，而不能简单合并。

在实现过程中，我遇到过一些很具体的错误。早期有些观测因为不能作为种子而被永久淘汰，后来发现这其实过于激进。一个观测不适合作为种子，不代表它不能作为某个已有假设的成员；如果过早淘汰，就会让候选缺少本来有用的视角证据。还有一个问题是“独立支持”的定义：如果两个观测只是共同参考了同一个已有 Mask3D 掩码，它们不一定真的互相支持。后来我把独立强支持和 Mask3D 参照辅助支持分开统计，避免把共同参照误认为多视角一致性。

另一个让我印象很深的问题是同帧关系的判断顺序。最初我更关注父子包含关系，也就是一个掩码是否包含另一个掩码。但后来发现，在同一帧中，如果两个区域存在明显深度硬冲突，那么应该优先判断冲突，而不是先看包含。否则，一个前景物体和背景区域可能因为二维包含关系被误判为同一实例的不同部分。这个细节让我认识到，几何关系的优先级非常重要，不能只看二维平面上的重叠。

点级证据图还帮助我把“候选失败”拆成更具体的原因。例如，有些候选失败是因为支持观测太少，有些是因为冲突边过多，有些是因为点级投票后剩余点太少，有些是因为已经被可靠 Mask3D 解释。把这些原因记录下来之后，我才知道下一步该修哪一类问题，而不是继续盲目调一个全局阈值。

但是点级证据图也暴露了新的问题。三维点本身没有明确的物体边界。用固定距离连通点云时，距离近不一定代表属于同一物体。例如门和墙、椅子腿和地面、柜子和旁边的物体，在点云中可能非常接近。点级候选容易出现背景污染、多物体错误合并、零碎候选和跨边界补全。也就是说，点级证据图可以帮助判断“不同视角看到的是不是同一个物体”，但它很难回答“这个物体的三维边界到底在哪里”。

举一个更直观的例子，门、白板、窗帘、电视、床垫这类平面或大片物体，经常和墙面、门框、支撑结构靠得很近。如果只从点级距离看，它们很容易被连成一个整体。再比如椅子腿、架子、细长边缘这类结构，本身就容易断裂成多个点级小块；如果连通半径太小，就会碎；如果连通半径太大，又会跨到别的物体。这说明点级方法面临一个两难：半径小会导致碎片，半径大又会导致污染。

这让我形成了第二个重要认识：一个中间诊断工具即使方向正确，也可能只解决了部分问题。点级证据图解决了多视角身份关系，却不能可靠表达物体几何边界。如果继续只调点级阈值，例如调距离、调覆盖率、调冲突比例，可能会陷入局部修补，而不是解决根本矛盾。正是因为这个认识，我后来才转向“超点前移”：让证据图负责判断不同观测是否属于同一物体，让超点图负责表达物体在三维空间中的几何范围和边界。

\section{案例三：超点前移与 largest\_cc 裁剪的经验}

为了解决物体边界问题，我后来把超点信息前移。这里的超点不是新的大模型，而是复用 Mask3D/ScanNet200 预处理阶段已有的几何超点。它的意义在于：点级证据图负责判断不同观测是否属于同一个物体，超点图负责表达三维空间中的几何区域和边界。

在这个阶段，我构建了几套超点候选进行对照：第一套是点级候选，第二套是 core-only，即只保留多视角强支持的核心超点，第三套是 core+boundary，即在核心超点之外补一层边界超点，第四套是在严格 core+boundary 后再保留最大点级连通分量，也就是 largest\_cc。这个诊断最有价值的地方在于，它不直接改变最终 AP，而是先比较候选点数、连通分量、点级覆盖、冲突重叠和 Mask3D 覆盖。

10 个场景的结果显示，largest\_cc 能显著减少碎片。点级候选只有 13/63 是单连通，严格 core+boundary 是 47/63，而 largest\_cc 达到 60/63。从表面上看，这似乎说明 largest\_cc 很成功。但是这时我没有直接把它接入 AP，而是做了人工可视化检查。这个步骤非常关键。

例如，在 scene0011\_00 的 dishwasher 候选中，largest\_cc 只裁掉了 2 个点，基本就是小碎片；在 scene0046\_02 的 toilet 候选中，裁掉了 14 个点，也像是边缘小团噪声。这说明 largest\_cc 对这些候选是安全的。但是在 door 候选中，虽然裁剪本身没有明显误删主体，但超点候选整体仍然比点级候选大很多，需要进一步判断它是在合理补全整扇门，还是把墙面、门框或相邻平面也补进来了。

这个案例给我的教训很深：连通性变好，不等于候选一定可用。largest\_cc 解决的是“碎片”问题，不解决“是否过度补全”问题。一个候选可以是单连通的，但仍然可能包含过多背景或相邻平面。如果当时只看 60/63 的单连通结果就直接跑 AP，很可能会误以为问题已经解决。

\section{案例四：从生成候选到判断候选动作}

在发现 largest\_cc 仍不能直接决定候选是否可用后，我进一步做了候选动作判断器的诊断版。它的目标不是再问“这个候选连不连通”，而是判断每个候选应该怎么处理：

\begin{itemize}
    \item \texttt{accept\_completion}：可以接受补全；
    \item \texttt{keep\_core\_only}：只保留核心，不要边界补全；
    \item \texttt{reject\_or\_needs\_mask3d\_support}：拒绝，或必须有 Mask3D 支持才用；
    \item \texttt{manual\_review}：需要人工复查。
\end{itemize}

这个判断器使用的特征包括：largest\_cc 与点级候选的点数比例、largest\_cc 被点级候选覆盖的比例、点级候选被 largest\_cc 覆盖的比例、Mask3D IoU 和覆盖率、冲突重叠比例，以及类别是否属于平面或大片类。平面/薄片类包括 whiteboard、tv、door、curtain、mattress、projector screen、bulletin board、mirror 和 mat；picture 虽然也是平面，但通常是小物体，所以单独标记，不能和 curtain、whiteboard 一样重罚。

这个阶段还暴露了一个规则设计上的问题：最早我把 Mask3D 支持定义得太宽，只要 seed coverage 高就认为有支持。但后来发现，这样会让一些候选仅仅因为落在已有 Mask3D 区域内部就被放行，即使它和 Mask3D 的 IoU 并不高。因此我把规则收紧为：Mask3D 支持不能只靠 seed coverage，必须有足够 IoU，或者 IoU 接近达标且 seed coverage 很高。

在 10 个场景上，v2 动作分布为：21 个 \texttt{accept\_completion}，32 个 \texttt{manual\_review}，9 个 \texttt{reject\_or\_needs\_mask3d\_support}，1 个 \texttt{keep\_core\_only}。几个重点例子也符合直觉：dishwasher 和 toilet 被判为可以补全；door 和 tv 被判为需要 Mask3D 支持或拒绝；picture 和 mini fridge 被判为人工复查；projector screen 在加入平面/薄片风险类后，也从接受改为人工复查。

这个案例说明，我的研究方式已经从“生成更多候选”转向“理解候选行为”。不同候选不应该被同等处理。小物体、紧凑物体、已有 Mask3D 支持的补全候选可以更积极；平面、大片、薄片类候选必须更保守；冲突重叠高或候选明显膨胀时，应该进入人工复查或只保留核心。

\section{总结：我对 Vibe Research 的理解变化}

回顾这次实践，我最大的体会是：Vibe Research 的价值不是让 AI 直接替我找到最终答案，而是帮助我更快地把模糊直觉变成可检查的实验和诊断。AI 确实极大提高了效率，例如帮助我快速写导出脚本、汇总 JSON/CSV、生成可视化图和审计规则。但是，如果研究者自己没有清晰的问题分解，AI 生成的代码和实验很可能只会扩大混乱。

这次经历中，我认为不正确的做法主要有三点。第一，只看最终 AP，不看中间证据。第二，看到某个指标变好就急于进入最终评估，例如 largest\_cc 单连通提升后差点忽略过度补全问题。第三，把规则设计得过宽，例如早期 Mask3D 支持只靠 seed coverage 放行，容易掩盖候选与真实边界不一致的问题。

相对而言，真正有效的做法是：先建立诊断，再决定是否进入 AP；先可视化重点反例，再扩大实验规模；先区分候选动作，再考虑融合策略。对我来说，这次 Vibe Research 最重要的产出不是某个单独模块，而是一套更可靠的研究流程：提出直觉、生成候选、建立诊断、发现反例、修正规则、再决定是否扩大实验。

如果用一句话概括我的经验，那就是：Vibe Research 可以很快，但不能很随意。真正有价值的 Vibe Research 不是“凭感觉乱试”，而是借助 AI 的速度，把研究者的直觉迅速转化为证据、反例和可复查的决策规则。

\end{document}
